
\section{Indeks potrójnych skrzyżowań}
\index{indeks!potrójnych skrzyżowań|(}%
% DICTIONARY;triple crossing number;indeks potrójnych skrzyżowań;-

Dotychczas rzutowaliśmy sploty tak, żeby w~każdym punkcie wielokrotnym spotykały się dokładnie dwa łuki.
Adams \cite{adams2013} zaproponował w~2013 roku, żeby zamiast tego pozwolić spotykać się tylko trzem łukom naraz, parami poprzecznie.
\index[persons]{Adams, Colin}%
Na diagramie oznacza się wtedy, który z~trzech łuków leży na górze, który pośrodku, a~który na dole.
Adams pokazał, że każdy splot ma taki diagram, a~zatem ma sens następująca definicja.

\begin{definition}
    Niech $L$ będzie splotem.
    Najmniejszą liczbę potrójnych skrzyżowań na diagramie, w~którym wszystkie skrzyżowania są potrójne, nazywamy indeksem potrójnych skrzyżowań splotu $L$ i~oznaczamy $\operatorname{c_3} L$.
\end{definition}

Potrójne skrzyżowanie można lekko zaburzyć, otrzymując trzy zwykłe skrzyżowania, dlatego $\crossing L \le 3 \operatorname{c_3} L$.
Dla splotów alternujących Adams udowodnił mocniejsze ograniczenie $\crossing L \le 2 \operatorname{c_3} L$.
Trójlistnik i~ósemka mają po dwa potrójne skrzyżowania.

Jabłonowski \cite{jablonowski2020} znalazł dolne ograniczenie przez genus kanoniczny: jeśli $L$ ma $\mu$ ogniw, to
\index[persons]{Jabłonowski, Michał}%
\index{genus!kanoniczny}%
\begin{equation}
    \operatorname{c_3} L \ge 2 \operatorname{g_c} L + \mu - 1.
\end{equation}
Nierówność okazuje się na tyle dobra, że wyznacza indeks potrójnych skrzyżowań wszystkich węzłów torusowych: $\operatorname{c_3} T_{p, q} = (p-1)(q-1)$.
\index{węzeł!torusowy}%
W~kolejnej pracy \cite{jablonowski2023} ten sam autor stablicował węzły pierwsze, których indeks potrójnych skrzyżowań nie przekracza pięciu.
Na przykład $5_2$ oraz $6_1$ mają ich trzy.

Adams wykorzystał diagramy z~potrójnymi skrzyżowaniami także do oszacowania objętości hiperbolicznej, wracamy do tego w~rozdziale siódmym.
\index{objętość!hiperboliczna}%

\index{indeks!potrójnych skrzyżowań|)}%

% Koniec podsekcji Indeks potrójnych skrzyżowań
